
\section{Introduction}
\label{intro}

% Designing FPGA-based accelerators for modern Artificial Intelligence (AI) workloads requires navigating a large and complex hardware design space that involves compute parallelism, memory hierarchies, dataflow strategies, and resource allocation constraints. As AI models continue to grow in complexity and deployment diversity, developing efficient accelerator architectures becomes increasingly time-consuming and heavily reliant on hardware expertise~\cite{mohaidat2024survey}.

Designing FPGA-based accelerators for modern artificial intelligence (AI) workloads requires selecting among interacting architectural choices, including compute parallelism, tiling factors, memory hierarchy organization, dataflow strategies, and on-chip resource allocation. These choices jointly affect latency, throughput, bandwidth pressure, and FPGA resource utilization, making accelerator design space exploration both expensive and expertise-intensive~\cite{mohaidat2024survey}. 
Although FPGA-based accelerators provide flexibility and energy efficiency compared to fixed-function ASICs, identifying performant accelerator configurations remains a significant challenge due to the large number of architectural permutations and optimization trade-offs.

Existing methodologies such as SECDA~\cite{haris2021secda} and toolkits like SECDA-TFLite~\cite{haris2023secda} enable rapid hardware software co-design through SystemC-based simulation and FPGA execution, significantly reducing the development overhead associated with custom accelerator design. They provide reusable architectural templates and automated hardware flows for AI acceleration. However, the process of exploring optimal accelerator configurations still requires extensive manual design space exploration (DSE), iterative tuning, and expert-driven reasoning over hardware trade-offs.

Recent advances in Large Language Models (LLMs) have demonstrated strong capabilities in reasoning, code generation, and structured problem solving across software engineering and systems domains~\cite{hou2024large, fan2023large}. These capabilities create an opportunity to assist or automate hardware accelerator exploration by enabling reasoning-guided design refinement and workload-aware architectural adaptation. However, applying LLMs to FPGA accelerator DSE remains relatively unexplored, particularly in combining LLM reasoning with structured hardware exploration flows and validating generated accelerator designs through real FPGA execution.

In our previous work~\cite{sharma2026llm}, we introduced \textbf{SECDA-DSE}, a framework that integrates LLM-guided reasoning into the SECDA ecosystem for FPGA accelerator DSE, and evaluated a single accelerator design using high-level simulation. 
SECDA-DSE combines a structured DSE Explorer, responsible for generating and evaluating candidate hardware accelerator configurations, with an LLM Stack that performs reasoning-guided refinement using Retrieval-Augmented Generation (RAG), Chain-of-Thought (CoT) prompting, and parameter-efficient fine-tuning. The framework incorporates a feedback-driven evaluation loop in which simulation and hardware execution metrics are collected and reused to iteratively improve subsequent decisions, as shown in Figure~\ref{fig:SECDA-DSE}.


In this paper, we present SECDA-DSE and extend its evaluation by generating multiple accelerator kernels and performing end-to-end real hardware execution on an FPGA. 
Specifically, we evaluate generated accelerators for element-wise vector multiplication, 2D-convolution, and matrix transpose. The generated designs successfully execute on FPGA hardware and exhibit workload-specific architectural characteristics and resource utilization patterns, highlighting the ability of the framework to adapt exploration decisions across different computational workloads without extensive fine-tuning.

This paper makes the following contributions:

\begin{itemize}
    \item We present SECDA-DSE, an LLM-guided framework for FPGA accelerator DSE that combines structured exploration, RAG, CoT prompting, fine-tuning, and iterative hardware evaluation.

    \item We extend the evaluation of SECDA-DSE from a single simulation-based accelerator design to multiple accelerator kernels validated through end-to-end FPGA execution, including vector multiplication, convolution, and transpose kernels. All generated accelerators passed functional validation on the FPGA platform, confirming their execution correctness. Representative execution results show measured latencies of 154 ms for vector multiplication, 163 ms for 2D convolution, and 238 ms for transpose. The generated designs also exhibit different resource profiles, such as 21.82\% DSP utilization for vector multiplication and 8.85\% LUT utilization for transpose.

    % \item We provide insights into workload-aware accelerator exploration and the potential of LLMs for enabling more automated hardware design workflows.
\end{itemize}